\documentclass[11pt,a4paper]{article}

\usepackage[T1]{fontenc}
\usepackage{lmodern}
\usepackage[margin=28mm]{geometry}
\usepackage{amsmath,amssymb,amsthm}
\usepackage{booktabs}
\usepackage{microtype}
\usepackage[hidelinks]{hyperref}

\newtheorem{theorem}{Theorem}
\newtheorem{lemma}[theorem]{Lemma}
\newtheorem{corollary}[theorem]{Corollary}
\newtheorem{proposition}[theorem]{Proposition}
\theoremstyle{remark}
\newtheorem{remark}[theorem]{Remark}
\newtheorem{example}[theorem]{Example}
\newcommand{\Ostar}{O^{*}}
\newcommand{\E}{\mathbb{E}}

\clubpenalty=10000
\widowpenalty=10000
\displaywidowpenalty=10000
\emergencystretch=2em
\hypersetup{pdftitle={The Key Lemma for Random Signs}}

\title{The Key Lemma for Random Signs}
\author{}
\date{8 October 2026\\[3pt]\small Working notes}

\begin{document}
\maketitle

\begin{abstract}
We prove the Key Lemma of the plan note~\cite{N7}. Consider the search
of~\cite{N7}: a fixed variable order, pruning clauses, and the sum check. Let
$\pi_t$ be the probability that a uniformly random partial assignment survives
depth~$t$. For any graph with independent uniform signs we bound
$\E\,\pi_T$ by four explicit terms, each exponentially small in $n$ when a
simple margin of the graph is of order $n\sqrt d$. Two consequences follow. On random
formulas $R(n,d)$ of every density $d_0\le d\le n-1$, the search solves Max
2-SAT in time $\Ostar(2^{(1-\gamma)n})$ for every $\gamma<0.0236$, with
probability $1-o(1)$. On every $d$-regular graph with random signs, in any
order, every $\gamma<0.0132$ works. The proof looks at the signs of a variable
only when the search reaches it. Then the three terms of the race identity
of~\cite{N7} are sums of fair coins that have not been looked at, and the
inequalities of Szarek, Azuma and McDiarmid finish the proof. The constant is
far below the measured saving of about $0.5n$. Random signs are used in exactly
three places, and these three places sharpen the worst-case target.
\end{abstract}

\section{Statement}\label{sec:statement}

We use the search of~\cite[Section~2]{N7}. The formula is canonical with
$B=1$, and the order $x_1,\dots,x_n$ is fixed in advance. Branching on $x_t$
gives the non-preferred child the pruning clause $\langle L_t\le p-1\rangle$
over the free neighbours of $x_t$. A node is cut when a clause is violated
(this includes $p=0$) or when the sum check fires. Let $\pi_t(F)$ be the
probability that a uniform $\sigma\in\{0,1\}^n$ survives the checks at depths
$1,\dots,t$; the plan note calls it $S_t$, but here $S$ denotes the shortfall.
By Knuth's identity~\cite[Proposition~1]{N7}, the search has at most
$1+2n\max_t2^t\pi_t$ nodes.

\paragraph{Graph notation.} $G$ is the graph of $F$, with one edge per clause.
For a variable $x_t$, $D^+_t$ counts its neighbours after it in the order (the
support of its pruning clause) and $D^-_t$ its neighbours before it. For a
depth $T$, $e_{\le T}=\sum_{t\le T}D^-_t$ is the number of edges among
$x_1,\dots,x_T$ and $e_T$ the number of edges between them and the rest. For a
free variable $x_w$, $w>T$, $\deg_T(w)$ is its number of neighbours among
$x_1,\dots,x_T$. \emph{Random signs} on $G$ means that each clause
$(\ell_u\vee\ell_v)$ gets independent uniform signs for $\ell_u$ and for
$\ell_v$. $R(n,d)$ is the random graph $G(n,d/(n-1))$ with random signs.

\begin{theorem}\label{thm:main}
Let $G$ be a graph on $n$ vertices, $1\le T\le n$, and put
\[
 \mu_G=\frac1{4\sqrt2}\sum_{t\le T}\sqrt{D^+_t+2D^-_t},\qquad
 U_G=\frac12\sum_{w>T}\sqrt{\deg_T(w)},\qquad
 V_G=\sum_{t\le T}\Bigl(\frac{A_t}8+\frac{\sqrt{A_t}}2+\frac12\Bigr),
\]
where $A_t=D^+_t+4D^-_t$. If $\alpha,\xi,\kappa,\upsilon\ge0$ and
$\alpha+\xi+\kappa+\upsilon<\mu_G-U_G$, then for random signs on $G$
\[
 \E\,\pi_T\le e^{-\alpha^2/(2V_G)}+e^{-16\xi^2/e_{\le T}}
 +e^{-2\kappa^2/e_{\le T}}+e^{-2\upsilon^2/e_T}.
\]
\end{theorem}

For $R(n,d)$ and for $d$-regular graphs, at a suitable depth $T$, the
margin $\mu_G-U_G$ is of order $n\sqrt d$ while $V_G$, $e_{\le T}$ and $e_T$
are $O(nd)$ (Section~\ref{sec:constants}). Each term is then
$e^{-\Omega(n)}$, independently of $d$. The theorem bounds a
single depth; that is enough.

\begin{lemma}[One depth suffices]\label{lem:onedepth}
For every $T$, $\max_t2^t\pi_t\le\max\{2^T,\,2^n\pi_T\}$.
\end{lemma}

\begin{proof}
$2^t\pi_t\le2^t\le2^T$ for $t\le T$, and $\pi_t\le\pi_T$ for $t>T$.
\end{proof}

\begin{corollary}[Key Lemma]\label{cor:key}
For every $\gamma<0.0236$ there is $d_0$ such that for all $n$ and all
$d_0\le d\le n-1$, with probability $1-o(1)$ over $F\in R(n,d)$,
\[
 \max_{0\le t\le n}2^t\pi_t(F)\le2^{(1-\gamma)n}.
\]
Hence the search solves Max 2-SAT on $R(n,d)$ in time $\Ostar(2^{(1-\gamma)n})$
and polynomial space, with probability $1-o(1)$, at every density.
\end{corollary}

\begin{corollary}[Regular graphs]\label{cor:regular}
For every $\gamma<0.0132$ there is $d_0$ such that for every $d$-regular graph
$G$ with $d\ge d_0$ and every order of its vertices, with probability
$1-e^{-\Omega(n)}$ over random signs on $G$,
$\max_t2^t\pi_t\le2^{(1-\gamma)n}$.
\end{corollary}

So the graph does not matter; only the signs do. This is the form we want for
the worst case (Section~\ref{sec:uses}).

\section{Looking at signs late}\label{sec:late}

The proof rests on one observation in clause terms: when the search reaches
$x_t$, it has never used the signs of $x_t$. We state it for the \emph{uncut
process}, which runs $\sigma$ through all steps, creates every clause and keeps
it. On the event that $\sigma$ survives depth $T$, the uncut process and the
search agree up to depth $T$.

For a pruning clause $c$ and a variable $x_w$ of its support, the \emph{bad
value} of $x_w$ in $c$ is the value that makes the literal of $x_w$ in $c$
true; $j_c$ in~\cite[Lemma~2]{N7} counts the support variables that took their
bad value. Concretely, let $c$ belong to $x_s$ with preferred literal $\ell$,
and let $(\ell_s\vee\ell_w)$ be the clause of $x_s$ and $x_w$. If
$\ell_s=\ell$, the literal of $x_w$ in $c$ is $\neg\ell_w$: since $\ell$ is
false in this child, the clause is falsified exactly when $\ell_w$ is false.
If $\ell_s=\neg\ell$, the literal is $\ell_w$, since a flip of $x_s$ back to
$\ell$ does not lose the clause if $\ell_w$ is true.

\begin{lemma}[Late signs]\label{lem:late}
Let $H_t$ consist of the graph, the signs of $x_1,\dots,x_t$ in all their
clauses, and $\sigma_1,\dots,\sigma_t$.
\begin{enumerate}
 \item[(a)] $H_t$ determines, for every $s\le t$, the effective gap $\Delta_s$,
 the preferred value of $x_s$, whether $x_s$ creates a clause, its allowance,
 and the values $S_s$ and $Q_s$.
 \item[(b)] Given $H_{t-1}$, the signs of $x_t$ and the value $\sigma_t$ are
 independent and uniform.
 \item[(c)] Let $c$ be the clause of $x_s$, and $x_w$ in its support. Given
 $H_s$, the bad value of $x_w$ in $c$ is a function of the sign of $x_w$ in
 its clause with $x_s$, and the two signs give the two bad values.
\end{enumerate}
\end{lemma}

\begin{proof}
(a) The effective gap of $x_s$ sums over the clauses of $x_s$. A clause with a
later variable is still a 2-clause and contributes $\pm1$, by the sign of
$x_s$. A clause with an earlier $x_u$ contributes $\pm2$, by the sign of
$x_s$, if the literal of $x_u$ is false (the clause is now a unit clause on
$x_s$), and $0$ if it is true. Whether it is true depends on the sign of $x_u$
and on $\sigma_u$. So $\Delta_s$ depends only on the signs of
$x_1,\dots,x_s$ and on $\sigma_1,\dots,\sigma_{s-1}$. The preferred value
follows. A clause is created if $\sigma_s$ is the other value, and its
allowance is $p_s-1=(D^+_s-|\Delta_s|)/2-1$. $S_s$ adds $|\Delta_u|/2+1$ over
the created clauses. $Q_s$ adds $[\sigma_u\text{ bad in }c]-\frac12$ over the
assigned support variables $x_u$ of the created clauses $c$, and by~(c) these
bad values depend on the signs of $x_u$, $u\le s$. (b) All signs and all
values are independent, and $H_{t-1}$ contains no sign of $x_t$. (c) This is
the case distinction above: $\ell$ and $\ell_s$ are determined by $H_s$, and
the literal of $x_w$ in $c$ is $\neg\ell_w$ or $\ell_w$.
\end{proof}

\begin{example}\label{ex:tiny}
Let $x_1$ have the clauses $(x_1\vee x_2)$, $(x_1\vee x_3)$,
$(x_1\vee\neg x_4)$ and $(\neg x_1\vee x_5)$, and no others. Then
$\Delta_1=1+1+1-1=2$, $D^+_1=4$, $p=(4-2)/2=1$, and the preferred literal is
$x_1$. The child $x_1=0$ receives
\[
 P_1=\langle\neg x_2+\neg x_3+x_4+x_5\le0\rangle ,
\]
with bad values $x_2=0$, $x_3=0$, $x_4=1$, $x_5=1$. Each bad value is read
off from the sign of $x_w$ in its clause with $x_1$. The uncut process needs
it only when $x_w$ is assigned, or at depth $T$ through $U_T$. When the search
reaches $x_2$, its clause with $x_1$ has become the unit clause $(x_2)$. It
adds $+2$ to $\Delta_2$, and the same positive sign of $x_2$ makes $x_2=0$ its
bad value in $P_1$. So the preferred value of $x_2$ tends to be good for
$P_1$. Such correlations exist between the three terms below, and we never
need independence between them: we add their failure probabilities. Had
$\sigma_1=1$, the clause $(x_1\vee x_2)$ would be satisfied and would add $0$
to $\Delta_2$.
\end{example}

In words, Lemma~\ref{lem:late} gives three sources of unseen fair coins.
\begin{enumerate}
 \item At step $t$, the gap $\Delta_t$ is $\pm1$ over the 2-clauses of $x_t$
 plus $\pm2$ over its unit clauses, with unseen signs.
 \item At step $t$, the bad values of $x_t$ in the pruning clauses that
 contain it are unseen fair coins.
 \item At depth $T$, the bad values of the free variables are unseen fair
 coins.
\end{enumerate}

\section{The three estimates}\label{sec:estimates}

By~\cite[Lemma~2, Remark~3]{N7}, applied to the set of all clauses created up
to depth $T$, a partial assignment that survives depth $T$ has
\[
 Z_T=S_T+Q_T-U_T\le0,
\]
with $S_T=\sum_c(|\Delta_c|/2+1)$, $Q_T=\sum_c(j_c-a_c/2)$ and
$U_T=\frac12\sum_{w>T}|b_w|$, all taken over the uncut process. The three
lemmas below bound the three terms.

\begin{lemma}[Shortfall]\label{lem:S}
$\Pr[S_T<\mu_G-\alpha-\xi]\le e^{-\alpha^2/(2V_G)}+e^{-16\xi^2/e_{\le T}}$.
\end{lemma}

\begin{proof}
Write $S_T=\sum_{t\le T}X_t$ with $X_t=[x_t\text{ creates a
clause}]\,(|\Delta_t|/2+1)\ge0$. Let $r_t$ be the number of unit clauses on
$x_t$, i.e.\ of earlier neighbours whose literal in the shared clause is
false. Given $H_{t-1}$, $\Delta_t$ is a sum of $D^+_t$ unseen signs with weight
$1$ and $r_t$ unseen signs with weight $2$. A clause is created exactly when
$\sigma_t$ is the non-preferred value, which happens with probability
$\frac12$ independently of $\Delta_t$. By Szarek's inequality~\cite{Sza76},
$\E|\sum_ia_i\varepsilon_i|\ge(\sum_ia_i^2)^{1/2}/\sqrt2$ for independent
uniform signs $\varepsilon_i$, and therefore
\[
 \E[X_t\mid H_{t-1}]=\tfrac12\bigl(\tfrac12\E[|\Delta_t|\mid H_{t-1}]+1\bigr)
 \ge\frac{\sqrt{D^+_t+4r_t}}{4\sqrt2}+\frac12=:\mu_t .
\]
Also $\E[X_t^2\mid H_{t-1}]=\frac12\E[(|\Delta_t|/2+1)^2\mid H_{t-1}]
\le\frac12(\sqrt{D^+_t+4r_t}/2+1)^2\le v_t$, where
$v_t=A_t/8+\sqrt{A_t}/2+\frac12$ uses $r_t\le D^-_t$.

\emph{Lower tail.} For $X\ge0$ and $\lambda\ge0$, $e^{-\lambda X}\le1-\lambda
X+\lambda^2X^2/2$, so $\E[e^{-\lambda X_t}\mid H_{t-1}]\le\exp(-\lambda
\mu_t+\lambda^2v_t/2)$. Hence $\exp(-\lambda\sum_{s\le t}(X_s-\mu_s)-\frac{\lambda^2}
2\sum_{s\le t}v_s)$ is a supermartingale. With $\lambda=\alpha/V_G$,
\[
 \Pr\Bigl[S_T\le\sum_{t\le T}\mu_t-\alpha\Bigr]\le e^{-\alpha^2/(2V_G)}.
\]
This is the standard one-sided bound for nonnegative increments~\cite{Mau03};
it needs no upper bound on $X_t$.

\emph{The sum of the $\mu_t$.} Given the graph, the events ``the literal of
$x_u$ in its clause with $x_t$ is false'', $u<t\le T$, are $e_{\le T}$
independent fair coins: the literal of $x_u$ is false when its sign disagrees
with $\sigma_u$, and the signs of $x_u$ are independent and uniform, also given
$\sigma_u$. $r_t$ counts the coins of the clauses
between $x_t$ and earlier variables, and changing one coin changes
$\sum_t\sqrt{D^+_t+4r_t}$ by at most $2$. By McDiarmid's
inequality~\cite{McD89}, this sum falls more than $x$ below its mean with
probability at most $e^{-x^2/(2e_{\le T})}$. Its mean is at least
$\sum_t(\sqrt{D^+_t+2D^-_t}-1)$: the elementary inequality
$\sqrt y\ge\sqrt\mu+\frac{y-\mu}{2\sqrt\mu}-\frac{(y-\mu)^2}{2\mu^{3/2}}$ for
$y\ge0$ gives $\E\sqrt Y\ge\sqrt{\E Y}-\mathrm{Var}(Y)/(2(\E Y)^{3/2})$, and
here $\mathrm{Var}(D^+_t+4r_t)=4D^-_t\le2\,\E Y$. The terms $\frac12$ in $\mu_t$
pay for the $-1$, because $\frac1{4\sqrt2}<\frac12$. So
$\sum_t\mu_t\ge\mu_G-x/(4\sqrt2)$ except with probability
$e^{-x^2/(2e_{\le T})}$. Taking $x=4\sqrt2\,\xi$ proves the lemma.
\end{proof}

\begin{lemma}[Martingale]\label{lem:Q}
$\Pr[Q_T<-\kappa]\le e^{-2\kappa^2/e_{\le T}}$.
\end{lemma}

\begin{proof}
At step $t$, $Q$ changes by $\sum_c([\sigma_t\text{ bad in }c]-\frac12)$ over
the $m_t\le D^-_t$ created clauses that contain $x_t$; $m_t$ is determined by
$H_{t-1}$. Given $H_{t-1}$, the bad values of $x_t$ in these clauses are
independent fair coins (Lemma~\ref{lem:late}(b),(c)), and $\sigma_t$ is
independent of them. So the indicators $[\sigma_t\text{ bad in }c]$ are
independent fair coins, and the step of $Q$ is $\frac12$ times a sum of $m_t$
independent signs. So $Q$ is a walk with fair steps $\pm\frac12$, at most
$e_{\le T}$ of them, and
$\E[e^{-\lambda(Q_t-Q_{t-1})}\mid H_{t-1}]=\cosh(\lambda/2)^{m_t}\le
e^{\lambda^2D^-_t/8}$. Azuma's argument~\cite{Azu67,Hoe63} with
$\lambda=4\kappa/e_{\le T}$ gives the claim.
\end{proof}

\begin{lemma}[Uncertainty]\label{lem:U}
$\Pr[U_T>U_G+\upsilon]\le e^{-2\upsilon^2/e_T}$.
\end{lemma}

\begin{proof}
Condition on $H_T$. A free variable $x_w$ lies in $m_w\le\deg_T(w)$ created
clauses, one for each neighbour among $x_1,\dots,x_T$ that created a clause,
and $b_w$ is the sum of $\pm1$ over them according to its bad values. By
Lemma~\ref{lem:late}(c) these are unseen fair coins, independent over all
pairs $(c,w)$. So $\E[|b_w|\mid H_T]\le\sqrt{m_w}\le\sqrt{\deg_T(w)}$ by
Jensen's inequality, and $\E[U_T\mid H_T]\le U_G$. $U_T$ depends on at most
$e_T$ unseen signs, and changing one of them changes $U_T$ by at most $1$.
McDiarmid's inequality~\cite{McD89} finishes the proof.
\end{proof}

\begin{proof}[Proof of Theorem~\ref{thm:main}]
If $\sigma$ survives depth $T$, the search and the uncut process agree up to
depth $T$, so $Z_T\le0$. Outside the three events of
Lemmas~\ref{lem:S}--\ref{lem:U},
\[
 Z_T\ge(\mu_G-\alpha-\xi)-\kappa-(U_G+\upsilon)>0 .
\]
So survival lies in the union of the three events, and averaging over the
signs gives $\E\,\pi_T\le$ the sum of the four failure probabilities.
\end{proof}

\section{The constants}\label{sec:constants}

\begin{proof}[Proof of Corollary~\ref{cor:key}]
Let $p=d/(n-1)$, $T=\tau n$, and $\eta>0$. In $G(n,p)$ the following hold with
probability $1-o(1)$ for $d\ge d_0(\eta)$.
\begin{itemize}
 \item $e_{\le T}\le(1+\eta)\tau^2nd/2$ and $e_T\le(1+\eta)\tau(1-\tau)nd$, by
 Chernoff bounds.
 \item $U_G\le\frac12\sqrt{(n-T)e_T}\le(1+\eta)\frac12(1-\tau)\sqrt\tau\,n\sqrt
 d$, by Cauchy--Schwarz.
 \item $V_G\le(1+\eta)(\tau+\frac32\tau^2)nd/8$, since
 $\sum_{t\le T}A_t=e_T+5e_{\le T}$ and $\sum_t\sqrt{A_t}\le\sqrt{T\sum_tA_t}$.
 \item $\mu_G\ge(1-\eta)\frac1{6\sqrt2}((1+\tau)^{3/2}-1)\,n\sqrt d$. The
 expected value of $D^+_t+2D^-_t$ is $(n+t-2)p$, so the mean of
 $\sum_t\sqrt{D^+_t+2D^-_t}$ is at least
 $\sum_t\sqrt{(n+t-2)p}-T$ by the inequality used in Lemma~\ref{lem:S}. One edge
 changes the sum by at most $1+\sqrt2$, so the Efron--Stein
 inequality~\cite{ES81} bounds its variance by $3.2\,nd$, and Chebyshev's
 inequality finishes.
\end{itemize}
Put $M(\tau)=\frac1{6\sqrt2}((1+\tau)^{3/2}-1)-\frac12(1-\tau)\sqrt\tau$,
measure $\alpha,\xi,\kappa,\upsilon$ in units of $n\sqrt d$, and let
$\eta\to0$. The four exponents of Theorem~\ref{thm:main} become $k_in$ times
the square of the corresponding parameter, with
\[
 k_1=\frac4{\tau+\frac32\tau^2},\quad k_2=\frac{32}{\tau^2},\quad
 k_3=\frac4{\tau^2},\quad k_4=\frac2{\tau(1-\tau)} .
\]
The best split of $M(\tau)$ makes all four exponents equal to
$c(\tau)=M(\tau)^2/(\sum_ik_i^{-1/2})^2$. At $\tau=0.976$ we get
$M=0.1976$ and $c=0.01638$, so
$\E\,\pi_T\le4e^{-(c-\varepsilon)n}$ for typical graphs. By Markov's inequality,
$\pi_T\le e^{-(c-2\varepsilon)n}$ except with probability $4e^{-\varepsilon n}$
over the signs. Lemma~\ref{lem:onedepth} gives
$\max_t2^t\pi_t\le2^{(1-\gamma)n}$ for every
$\gamma<\min\{1-\tau,\,c/\ln2\}=\min\{0.024,\,0.02364\}$.
\end{proof}

\begin{proof}[Proof of Corollary~\ref{cor:regular}]
In a $d$-regular graph $D^+_t+2D^-_t=d+D^-_t\ge d$ in every order, so
$\mu_G\ge T\sqrt d/(4\sqrt2)$. Moreover $U_G\le\frac12(n-T)\sqrt d$,
$V_G\le T(d/2+\sqrt d+\frac12)$, $e_{\le T}\le Td/2$ and $e_T\le(n-T)d$. The same
optimization, with $M(\tau)=\tau/(4\sqrt2)-\frac12(1-\tau)$ and
$k=(1/\tau,\,32/\tau,\,4/\tau,\,2/(1-\tau))$, gives $c=0.00915$ at
$\tau=0.986$ and $c/\ln2=0.0132$.
\end{proof}

\section{How far from the truth}\label{sec:truth}

\begin{table}[ht]
\centering
\begin{tabular}{lccccccc}
\toprule
 & $S_T$ & $\sum_t\mu_t$ & $\mu_G$ & $U_T$ & $U_G$ & $Q_T$ & $\E|\Delta_t|$ ratio\\
\midrule
$n=1000$, $d=100$ & 0.253 & 0.235 & 0.190 & 0.026 & 0.047 & $-0.001$ & 1.14\\
$n=2000$, $d=200$ & 0.249 & 0.222 & 0.191 & 0.027 & 0.048 & $-0.002$ & 1.14\\
$n=2000$, $d=1000$ & 0.238 & 0.205 & 0.191 & 0.027 & 0.048 & $+0.001$ & 1.13\\
$n=1500$, complete & 0.224 & 0.202 & 0.191 & 0.028 & 0.047 & $-0.002$ & 1.14\\
\bottomrule
\end{tabular}
\caption{The terms at depth $T=0.9n$ in units of $n\sqrt d$, averaged over
four assignments; the bounds of the proof are $\mu_G$ and $U_G$. The last
column is the average of $|\Delta_t|/\sqrt{(D^+_t+4r_t)/2}$, against $1$ in
Szarek's inequality.}\label{tab:terms}
\end{table}

Table~\ref{tab:terms} compares the proof with the process. Every bound holds
with room, and the race is won by a wide margin: $Z_T$ is between
$0.19\,n\sqrt d$ and $0.23\,n\sqrt d$, where the proof uses
$\mu_G-U_G=0.143\,n\sqrt d$. The constant
$\gamma\approx0.02$ is small for three reasons.
\begin{enumerate}
 \item Four separate tails, each with a second-moment rate, are added.
 \item Szarek's constant $1/\sqrt2$ is used where the measured ratio is
 $\sqrt{4/\pi}\approx1.13$ times larger.
 \item Only one depth is used. The proof places the cut before depth
 $0.976n$, while the measured median cut is at $0.43n$ and the measured saving
 is $0.50n$ to $0.62n$~\cite[Tables~1 and~2]{N7}.
\end{enumerate}
A single exponential moment of $Z_t$ at a depth near $0.5n$ would remove the
first and the third loss; the Gaussian heuristic of~\cite[Section~6]{N7}
suggests a saving of about $0.2n$. We leave this; the point here is that the
saving is independent of the density.

\section{What the proof uses}\label{sec:uses}

Random signs enter only through Lemma~\ref{lem:late}, in three places.
\begin{enumerate}
 \item[(S)] The signs of $x_t$ in its 2-clauses and unit clauses make
 $|\Delta_t|$ large and the shortfall concentrated.
 \item[(Q)] The bad values of $x_t$ in the clauses that contain it make $Q$ a
 walk with fair steps $\pm\frac12$.
 \item[(U)] The bad values of the free variables make $|b_w|$ of order
 $\sqrt{m_w}$ instead of $m_w$.
\end{enumerate}
For a fixed formula, randomness comes only from $\sigma$. Two things remain. $Q$
is still a martingale, but its step at $x_t$ is $\pm|b_{x_t}|/2$. The shortfall
still has mean at least $\sum_{t\le T}(\sqrt{D^-_t/2}/4+\frac12)$, because each
earlier neighbour of $x_t$ satisfies or falsifies its literal by an independent
fair coin, so $\Delta_t$ is a fixed number plus $\pm1$ over the $D^-_t$ earlier
neighbours. What is lost is the size of the bad-value sums. The proof of
Theorem~\ref{thm:main} goes through verbatim for a fixed formula along every
path on which
\begin{enumerate}
 \item[(W1)] $U_T$ is a small constant times $n\sqrt d$,
 \item[(W2)] $\sum_{t\le T}b_{x_t}^2=O(nd)$, and
 \item[(W3)] $\sum_{t\le T}|\Delta_t|$ is at least a constant times $n\sqrt d$
 and $\sum_{t\le T}\Delta_t^2=O(nd)$,
\end{enumerate}
provided these hold except with probability $e^{-\Omega(n)}$ over $\sigma$.
(W1) and (W2) are cancellation of bad values, at the free and at the assigned
variables. (W3) says that effective gaps keep their typical size $\sqrt d$; it
makes the shortfall concentrate by the argument of Lemma~\ref{lem:S}. This
replaces the single cancellation condition of~\cite[Section~7]{N7}, which named
only (W1); the claim there that the shortfall part holds for every formula is
true only for its mean. Where (W1) or (W2) fails, many clauses agree about the
values of many variables. That is the aligned structure of~\cite[Section~9]{N4},
which is the case for aggregation.

\section{Checks}\label{sec:checks}

\texttt{explore\_key\_lemma\_proof.py} (a few seconds each):
\begin{itemize}
 \item \texttt{inclusion}: on random formulas with $n\le400$, all $51{,}041$
 pairs $(\sigma,t)$ in which $\sigma$ survives depth $t$ have $Z_t\le0$ for the
 uncut process, and the identity of~\cite[Lemma~2]{N7} holds exactly at every
 step.
 \item \texttt{terms}: Table~\ref{tab:terms}.
 \item \texttt{constant}: the optimizations in the proofs of
 Corollaries~\ref{cor:key} and~\ref{cor:regular}.
\end{itemize}

\begin{thebibliography}{99}
\raggedright
\bibitem{N4} Max 2-SAT with pruning clauses: pair decisions, sums, and a
 measure. Working note, 6--7 October 2026
 (\texttt{max2sat\_pairs\_sums\_measure.tex}).
\bibitem{N7} A simplified plan for Max 2-SAT: one search, one count, one Key
 Lemma. Working note, 8 October 2026 (\texttt{max2sat\_plan\_key\_lemma.tex}).
\bibitem{Azu67} K.~Azuma. Weighted sums of certain dependent random variables.
 \emph{T\^ohoku Mathematical Journal} 19(3):357--367, 1967.
\bibitem{ES81} B.~Efron and C.~Stein. The jackknife estimate of variance.
 \emph{The Annals of Statistics} 9(3):586--596, 1981.
\bibitem{Hoe63} W.~Hoeffding. Probability inequalities for sums of bounded
 random variables. \emph{Journal of the American Statistical Association}
 58(301):13--30, 1963.
\bibitem{Mau03} A.~Maurer. A bound on the deviation probability for sums of
 non-negative random variables. \emph{Journal of Inequalities in Pure and
 Applied Mathematics} 4(1), Article~15, 2003.
\bibitem{McD89} C.~McDiarmid. On the method of bounded differences. In
 \emph{Surveys in Combinatorics 1989}, London Mathematical Society Lecture Note
 Series 141, pages 148--188. Cambridge University Press, 1989.
\bibitem{Sza76} S.~J.~Szarek. On the best constants in the Khinchin inequality.
 \emph{Studia Mathematica} 58(2):197--208, 1976.
\end{thebibliography}

\end{document}
